\section{Testavgrensninger og åpne hull}
\label{sec:testavgrensninger}

Testene dekker funksjonelle krav, sentrale sikkerhetskrav og validering av
provider-arkitekturen, men de dokumenterer ikke produksjonsklarhet.

Det ble ikke gjennomført automatisk sikkerhetsskanning i CI, lasttesting med
dedikerte verktøy, systematisk fuzz-testing, kontraktstesting mot flere versjoner
av AntMedia/NHN/Pexip eller automatiserte recovery-tester for container-restart
og nettverksfeil. TLS-konfigurasjon er støttet, men testsuiten verifiserer ikke
sertifikatkjede, TLS-handshake eller sertifikatrotasjon.

E2E-skriptet gir funksjonell validering over tid, men lagrer ikke p95/p99,
gjennomstrømming eller maksimalt antall samtidige rom som målbare resultater.
Derfor kan testene brukes som bevis for at proof-of-concept-løsningen fungerer i
testmiljøet, men ikke som bevis for produksjonsmessig robusthet.

Avgrensningene svekker ikke proof-of-conceptets verdi, men markerer hva
som må adresseres før en eventuell produksjonssetting med reelle
pasientdata og strengere driftskrav.
